\section{Limitations}
\label{sec:limitations}

\sys{} has several concrete limitations, some of which are artefacts of
this study and some of which are structural.

\paragraph{Small evaluation set.}
POOL-A contains four competition runs and POOL-B contains tens of
single-trick segments. This is a case study, not a statistically powered
benchmark. We report per-clip numbers rather than only aggregate means so
readers can form their own picture, and we release the clips and labels
so the community can extend the set. A larger, multi-annotator benchmark
is the natural next paper.

\paragraph{Single-annotator ground truth.}
All labels were assigned by the author. For tricks with clear visual
identity (backflip, frontflip, gainer), this is uncontroversial. For
multi-twist tricks and off-axis variants, human judges themselves
disagree, and a single annotator is insufficient. We release
\texttt{ground\_truth.csv} with a \texttt{notes} column so external
annotators can propose corrections.

\paragraph{Closed-model dependence.}
The strongest configurations in the leaderboard are closed-weight
frontier models (Gemini 2.5, Claude Sonnet 4.5). Their behaviour may
change between versions, which is a reproducibility risk. We mitigate
this by (i) always including an open-weight comparison point
(Qwen2.5-VL 72B), (ii) logging the exact model version string in every
\texttt{results.jsonl} record, and (iii) releasing the full prompt and
fuzzy matcher so the \emph{system} is reproducible even if specific
weights drift.

\paragraph{Monocular video only.}
Real competitions film from multiple synchronised angles. A
multi-camera extension of \sys{} would plausibly resolve most
off-axis and wall-context failures documented in
Sec.~\ref{sec:experiments:failures}, but is out of scope here. We
note only that the fuzzy matcher and D-score assembler are camera-count
agnostic and could be reused as-is.

\paragraph{D-score only.}
FIG judging decomposes into a D-score (difficulty, from the table of
tricks) and an E-score (execution, deductions for safety, landing
quality, and flow). \sys{} addresses only the D-score; landing quality
and flow scoring require frame-level body-pose analysis that VLMs are
currently not precise enough to provide. We explicitly position the
E-score as future work, not as a broken promise.

\paragraph{Segmentation robustness.}
The inversion segmenter assumes the athlete becomes inverted during the
aerial phase. This holds for all acrobatic and wall tricks in the FIG
2025 table, but misses low-rotation moves (swing 180, wall spins) and
is confused by crowd members who happen to bend over at the wrong
moment. A context-aware segmenter driven by an additional VLM call is a
natural next step.
